%=============================================================================
\section*{两个支配世界的数：$e$ 与 $\pi$}
\addcontentsline{toc}{section}{两个支配世界的数：$e$ 与 $\pi$}
%=============================================================================

在正式进入滤波器、传递函数和各种控制工具之前，让我们先花一点时间认识两个无理数——$\pi$ 和 $e$。它们悄无声息地支撑着这本笔记里几乎所有的公式。

你在高数课上见过它们，可能还背过小数点后好几位。但下面一句话就能概括\textit{它们在工程中为什么重要}：

\begin{mdframed}
\textbf{$\pi$ 是"重复"的语言。}\\
\textbf{$e$ 是"增长与衰减"的语言。}
\end{mdframed}

就这么简单。这本笔记里你将遇到的每一个方程，几乎都建立在这两个思想之上——或者同时建立在它们两个之上。

%---------------------------------------------------------------------
\subsection*{$\pi$：周期性的心跳}
%---------------------------------------------------------------------

转一个轮子，看一个钟摆来回摆动，听一根音叉嗡嗡作响——这些全是\textit{周期性}现象：同样的轨迹不断重复。而 $\pi$ 就是把圆的几何和重复的代数联系在一起的那个数。

想象一个点以匀速在单位圆上运动。它在 $t$ 时刻的水平坐标是 $\cos(t)$，垂直坐标是 $\sin(t)$。绕一整圈恰好需要 $2\pi$ 的时间。这不是巧合——这正是 $\pi$ 的来源。圆的周长是 $2\pi r$；除以半径就得到 $2\pi$，即旋转的"自然周期"。

在控制工程里，周期现象无处不在：机械结构的振动、电源的交流电压、欠阻尼阶跃响应中的振荡（系统超调后来回振荡、最终稳定下来的过程——后面会详细讲）、单片机的采样时钟——全都是周期性的。

\begin{figure}[H]
  \centering
  \includegraphics[width=\textwidth]{circle_sinusoid.pdf}
  \caption{单位圆上的一个点追踪出 $\cos t$ 与 $\sin t$。绕一整圈恰好 $2\pi$。}
\end{figure}

每当你写下 $\sin(\omega t)$ 或 $\cos(\omega t)$，$\pi$ 就藏在 $\omega = 2\pi f$ 里。傅里叶变换——我们分析信号的核心工具——把\textit{任意}信号分解成一系列正弦波的叠加，每一个正弦波都带着 $\pi$ 的因子。所以 $\pi$ 不只是关于圆的；它是"这个东西多久重复一次？"这个问题的通用语言。

%---------------------------------------------------------------------
\subsection*{$e$：增长与衰减的标尺}
%---------------------------------------------------------------------

现在换一个完全不同的场景：一个电容通过电阻放电。电压不会振荡——它只是慢慢减小。你每秒测一次，发现它总是损失\textit{剩余值的同一比例}。放射性衰变中的半衰期也是一样的道理：真正重要的不是绝对减少了多少，而是\textit{剩下了多大比例}。

这种"恒定比例变化"正是 $e$ 所刻画的。形式上，$e$ 是唯一满足 $f(x) = e^x$ 等于自身导数的底数：
\[
  \frac{d}{dx}\,e^{x} = e^{x}.
\]
用大白话说：$e^x$ 在任意一点的变化速度等于它当前的值。这使得 $e^x$ 成为"变化速度与当前状态成正比"这一方程的天然解——而这正是微分方程 $\dot{x} = ax$ 所描述的。

当 $a > 0$ 时，$x(t) = e^{at}$ 指数增长——像复利，像失稳系统一路狂飙。当 $a < 0$ 时，$x(t) = e^{at}$ 指数衰减——像那个正在放电的电容，像稳定系统慢慢回到平衡点。

\begin{figure}[H]
  \centering
  \includegraphics[width=0.75\textwidth]{exponential_growth_decay.pdf}
  \caption{指数曲线 $e^{at}$。$a>0$：不稳定增长；$a<0$：稳定衰减。}
\end{figure}

在控制理论中，$a$ 的正负（更精确地说，系统极点的实部）决定了稳定性的一切。极点在左半平面意味着 $e^{at}$（$a<0$）：响应会衰减。极点在右半平面意味着 $a>0$：响应会发散。就这一个想法——指数是在变大还是在变小——就是稳定性分析的全部基础。

到这里，$e$ 管衰减，$\pi$ 管振荡，两者泾渭分明。但真正让人拍案叫绝的是——它们其实可以合二为一。

%---------------------------------------------------------------------
\subsection*{当它们联手：$e^{j\theta}$}
%---------------------------------------------------------------------

真正的魔法发生在 $e$ 和 $\pi$ 相遇的时候。欧拉公式把它们绑在了一起：
\[
  e^{j\theta} = \cos\theta + j\sin\theta,
\]
其中 $j = \sqrt{-1}$。这个等式可以说是整个工程数学中最重要的恒等式。它告诉我们：\textit{复指数}同时既是指数函数又是正弦函数——增长/衰减和振荡，统一在一个表达式里。

当系统具有复数极点 $\sigma \pm j\omega$ 时，时域响应为
\[
  e^{\sigma t}\bigl(\cos\omega t + j\sin\omega t\bigr).
\]
其中 $e^{\sigma t}$（$e$ 的部分）控制\textbf{包络}——振荡是在增大、衰减还是保持不变；$\cos\omega t$（$\pi$ 的部分，因为 $\omega = 2\pi f$）控制\textbf{频率}——系统振荡得有多快。

\begin{figure}[H]
  \centering
  \includegraphics[width=0.82\textwidth]{damped_oscillation.pdf}
  \caption{阻尼振荡 $e^{\sigma t}\cos(\omega t)$：包络在衰减（$e$），信号在振荡（$\pi$）。}
\end{figure}

换句话说：\textbf{$e$ 告诉你衰减多快，$\pi$ 告诉你振荡多频}。这本笔记里你将研究的每一个二阶系统——从质量-弹簧-阻尼器到 PID 控制的电机——都不过是这两个数各司其职而已。

\vspace{1em}
\noindent $\pi$ 和 $e$ 已经就位，接下来的一切都有了根基。
